\documentclass[12pt, a4paper]{article}

\usepackage[utf8]{inputenc}
\usepackage[T2A]{fontenc}
\usepackage[russian]{babel}
\usepackage{amsmath, amssymb, amsthm}
\usepackage{geometry}
\usepackage{enumitem}
\usepackage{graphicx}
\usepackage{float}
\usepackage{xcolor}

\geometry{top=2cm, bottom=2cm, left=3cm, right=1.5cm}

\newtheorem{theorem}{Теорема}
\renewcommand{\qedsymbol}{$\blacksquare$}

\title{Билет №76 \\ \small Обобщенная архитектура, состав и функции системы управления базой данных (СУБД). Характеристика современных технологий БД. Примеры соответствующих СУБД. (СпБГУ)}
\author{Сериков Владислав Александрович}
\date{16 мая 2026}

\begin{document}

\maketitle
\tableofcontents
\newpage

\section{Введение. Базовые определения}

Для обеспечения связности материала введем базовые понятия.

\textbf{База данных (БД)} --- это организованная структура, предназначенная для хранения, изменения и обработки взаимосвязанной информации, преимущественно больших объемов.

\textbf{Система управления базами данных (СУБД)} --- это комплекс программных и языковых средств, необходимых для создания баз данных, поддержания их в актуальном состоянии и организации поиска в них необходимой информации.

\section{Обобщенная архитектура СУБД}

В большинстве современных систем используется трехуровневая архитектура \textbf{ANSI/SPARC}, предложенная в 1975 году. Цель данной архитектуры --- отделение пользовательского представления данных от их физического хранения (обеспечение независимости данных).

\begin{enumerate}
    \item \textbf{Внешний уровень (Уровень представлений).} Самый высокий уровень, с которым работают конечные пользователи или приложения. Описывает только ту часть БД, которая релевантна для конкретного пользователя, скрывая остальные данные (представления/views).
    \item \textbf{Концептуальный (логический) уровень.} Описывает, \textit{какие} данные хранятся в БД и какие связи существуют между ними. Это полное логическое описание всей БД, не зависящее от физического хранения.
    \item \textbf{Внутренний (физический) уровень.} Наиболее низкий уровень архитектуры. Описывает, \textit{как} именно данные хранятся на физических носителях (файловые структуры, индексы, B-деревья, хеш-таблицы).
\end{enumerate}

\section{Состав и основные функции СУБД}

Обобщенно СУБД состоит из следующих основных компонентов (подсистем):

\begin{itemize}
    \item \textbf{Процессор запросов (Query Processor):} включает в себя лексический и синтаксический анализаторы, оптимизатор запросов и механизм выполнения запросов. Переводит высокоуровневые запросы (например, SQL) в низкоуровневые инструкции.
    \item \textbf{Менеджер хранения (Storage Manager):} управляет распределением пространства на диске, включает менеджер буфера (отвечает за перенос блоков данных из оперативной памяти на диск и обратно).
    \item \textbf{Менеджер транзакций (Transaction Manager):} обеспечивает выполнение свойств ACID (Atomicity, Consistency, Isolation, Durability) для транзакций, управляет блокировками для обеспечения конкурентного доступа.
    \item \textbf{Система журналирования (Log Manager):} ведет журнал упреждающей записи (WAL --- Write-Ahead Logging) для восстановления данных после сбоев.
\end{itemize}

\textbf{Основные функции СУБД:}
1. Управление данными во внешней памяти.
2. Управление буферами оперативной памяти.
3. Управление транзакциями.
4. Журналирование изменений и восстановление базы данных после сбоев.
5. Поддержка языков БД (DDL, DML).
6. Управление доступом и обеспечение безопасности.

\section{Алгоритмы в СУБД: Поиск по B-дереву}

Большинство реляционных СУБД используют B-деревья (или B+-деревья) для организации индексов на внутреннем уровне архитектуры. Рассмотрим алгоритм поиска ключа в классическом B-дереве.

\subsection*{Алгоритм: Поиск ключа $K$ в B-дереве}

\textbf{Вход:} Корень B-дерева $x$, искомый ключ $K$. \\
\textbf{Выход:} Узел, содержащий ключ $K$ и его индекс в узле, либо \texttt{NULL}, если ключ не найден.

\textbf{Шаги алгоритма:}
\begin{enumerate}[label=\arabic*.]
    \item Инициализировать индекс $i = 1$.
    \item Пока $i \le x.n$ (где $x.n$ — количество ключей в узле $x$) и $K > x.key_i$, увеличивать $i$ на 1 (линейный или бинарный поиск внутри узла).
    \item Если $i \le x.n$ и $K = x.key_i$, вернуть $(x, i)$ --- ключ найден.
    \item Если узел $x$ является листовым ($x.leaf = true$), вернуть \texttt{NULL} (поиск завершен безуспешно, так как дошли до листа, но не нашли ключ).
    \item Иначе прочитать с диска дочерний узел $x.child_i$ (выполнить \texttt{Disk-Read(x.child\_i)}).
    \item Рекурсивно или итеративно повторить алгоритм для узла $x.child_i$ и ключа $K$.
\end{enumerate}

\subsection*{Минимальный пример выполнения}
Рассмотрим B-дерево. \\
Корень содержит ключи: $[20, 40]$.\\
Дочерние узлы: $C_1 = [10, 15]$, $C_2 = [25, 30]$, $C_3 = [45, 50]$.

\textbf{Задача:} Найти ключ $K = 25$.
\begin{itemize}
    \item \textbf{Шаг 1-2 (Корень):} Сравниваем $K=25$ с ключами корня. $25 > 20$, но $25 < 40$. Следовательно, ключ лежит между 20 и 40.
    \item \textbf{Шаг 4-5:} Узел не листовой. Переходим к указателю $C_2$, соответствующему интервалу $(20, 40)$. Читаем $C_2$ с диска.
    \item \textbf{Шаг 1-3 (Узел $C_2$):} Сравниваем $K=25$ с ключами узла $C_2$. Первый же ключ $x.key_1 = 25$. $25 = 25$.
    \item \textbf{Результат:} Возвращаем указатель на запись, связанную с ключом 25.
\end{itemize}

\section{Теоретические основы современных СУБД: Теорема CAP}

С появлением распределенных СУБД фундаментальную роль приобрела теорема CAP, описывающая ограничения при проектировании таких систем.

\begin{theorem}[Теорема CAP / Теорема Брюера]
В любой распределенной системе хранения данных невозможно одновременно обеспечить выполнение более чем двух из следующих трех свойств:
\begin{itemize}
    \item \textbf{Consistency (Согласованность):} Любое чтение возвращает самую последнюю записанную версию данных или ошибку. (В строгом смысле — строгая линеаризуемость).
    \item \textbf{Availability (Доступность):} Любой запрос к не упавшему узлу завершается корректным, не содержащим ошибки ответом (однако без гарантии, что данные самые свежие).
    \item \textbf{Partition tolerance (Устойчивость к разделению):} Система продолжает функционировать несмотря на потерю или задержку произвольного числа сообщений между узлами сети.
\end{itemize}
\end{theorem}

\begin{proof}
Доказательство проведем от противного. Предположим, существует распределенная система, которая одновременно удовлетворяет свойствам C, A и P.

Пусть система состоит из двух узлов: $N_1$ и $N_2$.
\begin{enumerate}
    \item Так как система обладает свойством \textbf{P} (Partition tolerance), предположим, что между $N_1$ и $N_2$ произошел обрыв связи (network partition). Узлы работают, но не могут обмениваться сообщениями друг с другом.
    \item Клиент отправляет запрос на запись нового значения переменной $v = 1$ (до этого было $v = 0$) на узел $N_1$.
    \item Узел $N_1$ успешно сохраняет $v = 1$. Однако из-за сетевого разделения он не может передать обновление на узел $N_2$.
    \item Другой клиент отправляет запрос на чтение значения переменной $v$ на узел $N_2$.
    \item Так как система обладает свойством \textbf{A} (Availability), узел $N_2$ обязан ответить клиенту, не дожидаясь восстановления связи с $N_1$.
    \item Однако $N_2$ не получил обновление от $N_1$. Следовательно, узел $N_2$ вынужден вернуть старое значение $v = 0$.
    \item Возврат старого значения $v = 0$ нарушает свойство \textbf{C} (Consistency), требующее возврата последнего записанного значения ($v = 1$).
\end{enumerate}
Полученное противоречие доказывает, что невозможно обеспечить C, A и P одновременно в случае сетевого разделения. Следовательно, разработчик распределенной СУБД должен выбирать между связками CP (Согласованность + Устойчивость) и AP (Доступность + Устойчивость). (Системы CA возможны только в отсутствии разделения сети, что в реальных распределенных системах недостижимо).
\end{proof}

\section{Характеристика современных технологий БД}

Современные приложения требуют обработки колоссальных объемов данных (Big Data), высоких скоростей и гибких схем. Это привело к эволюции СУБД и появлению новых архитектур.

\subsection{Реляционные СУБД (RDBMS)}
Основаны на реляционной алгебре. Данные хранятся в таблицах со строгой схемой. Поддерживают SQL и строгие гарантии ACID. Вертикально масштабируются хорошо, горизонтально — с трудом (требуется шардирование).

\subsection{Технологии NoSQL (Not Only SQL)}
Созданы для горизонтального масштабирования, обработки неструктурированных данных и высокой производительности. Обычно жертвуют строгой согласованностью (выбирая AP по теореме CAP и модель BASE — Basically Available, Soft state, Eventual consistency).
Разделяются на 4 основные категории:
\begin{itemize}
    \item \textbf{Ключ-значение (Key-Value):} Данные хранятся как пары «ключ-значение». Максимально быстры для простых выборок.
    \item \textbf{Документоориентированные:} Данные хранятся в виде документов (JSON, BSON). Схема динамическая.
    \item \textbf{Колоночные (Wide-column):} Данные хранятся в виде разреженных матриц, сгруппированных по семействам колонок. Идеальны для аналитики и временных рядов.
    \item \textbf{Графовые:} Данные представлены в виде узлов, ребер и их свойств. Оптимизированы для сложных связей (социальные сети, рекомендательные системы).
\end{itemize}

\subsection{Технологии NewSQL}
Класс современных реляционных СУБД, которые стремятся объединить гарантии ACID и поддержку SQL (как в RDBMS) с горизонтальной масштабируемостью (как в NoSQL). Для достижения этой цели используются консенсусные алгоритмы (Raft, Paxos) и глобальные часы (например, TrueTime в Google Spanner).

\subsection{In-Memory базы данных}
СУБД, которые хранят все рабочие данные в оперативной памяти (RAM), а не на дисковых накопителях. Диск используется только для асинхронного журналирования транзакций (snapshot/WAL). Обеспечивают минимальную задержку (микросекунды).

\section{Примеры соответствующих СУБД}

В зависимости от применяемых технологий выделяют следующих лидеров индустрии:

\begin{itemize}
    \item \textbf{RDBMS:} PostgreSQL, Oracle Database, MySQL, Microsoft SQL Server.
    \item \textbf{NoSQL (Ключ-значение):} Amazon DynamoDB, Riak.
    \item \textbf{NoSQL (Документные):} MongoDB, Couchbase.
    \item \textbf{NoSQL (Колоночные):} Apache Cassandra, HBase, ClickHouse (аналитическая колоночная СУБД).
    \item \textbf{NoSQL (Графовые):} Neo4j, Amazon Neptune.
    \item \textbf{NewSQL:} CockroachDB, Google Spanner, YugabyteDB.
    \item \textbf{In-Memory:} Redis, Memcached, SAP HANA.
\end{itemize}

\end{document}
